\title{Chain-of-Thought Prompting Elicits Reasoning in Large Language Models}

\begin{abstract}
We investigate how producing a \textit{chain of thought}—a sequence of intermediate reasoning steps—can greatly enhance the capability of large language models to handle complex reasoning tasks. Specifically, we demonstrate that these reasoning skills emerge organically in sufficiently large language models through a straightforward technique called \textit{chain-of-thought prompting}, where a small number of chain-of-thought examples are included as prompts.

Experiments across three large language models reveal that chain-of-thought prompting boosts performance on various arithmetic, commonsense, and symbolic reasoning challenges. The improvements can be substantial. For example, providing PaLM 540B with just eight chain-of-thought exemplars results in state-of-the-art accuracy on the GSM8K benchmark for math word problems, outperforming even finetuned GPT-3 combined with a verifier.
\end{abstract}

\section{Introduction}
Recent advances in NLP have been driven by language models \citep[][\textit{inter alia}]{peters-etal-2018-deep,devlin-etal-2019-bert,brown2020language}. Increasing the size of these models has been shown to bring a variety of benefits, including enhanced performance and improved sample efficiency \citep[\textit{inter alia}]{kaplan2020scaling,brown2020language}. Nevertheless, merely enlarging the model scale has not sufficed to achieve strong results on difficult tasks such as arithmetic, commonsense, and symbolic reasoning \citep{rae2021scaling}.

In this study, we examine how a simple approach can unlock the reasoning potential of large language models, inspired by two concepts. First, methods for arithmetic reasoning gain from generating natural language explanations that lead to the final answer. Previous research has enabled models to produce such intermediate natural language steps either by training from scratch \citep{ling-etal-2017-program} or by finetuning pretrained models \citep{cobbe2021training}, as well as through neuro-symbolic techniques that employ formal languages rather than natural language \citep{roy-roth-2015-solving, chiang-chen-2019-semantically,amini-etal-2019-mathqa, chen2019neural}. Second, large language models provide the promising ability to perform few-shot in-context learning through \emph{prompting}. Rather than finetuning a dedicated language model checkpoint for each new task, one can simply prompt the model with a handful of input-output examples illustrating the task. Notably, this approach has proven effective for a range of straightforward question-answering tasks \citep{brown2020language}.

Both concepts mentioned above, however, have significant drawbacks. For rationale-augmented training and finetuning techniques, generating a substantial collection of high-quality rationales is expensive and far more complex than using simple input--output pairs typical in standard machine learning. Meanwhile, the conventional few-shot prompting approach employed in \citet{brown2020language} tends to perform poorly on tasks requiring reasoning skills and often shows limited improvement as the language model size increases \citep{rae2021scaling}. In this study, we integrate the advantages of these two approaches while circumventing their respective shortcomings. Concretely, we investigate the capacity of language models to conduct few-shot prompting for reasoning tasks using a prompt that contains triples: $\langle$input, \emph{chain of thought}, output$\rangle$.

A \emph{chain of thought} consists of a sequence of intermediate natural language reasoning steps that culminate in the final answer, and we designate this technique as \textit{chain-of-thought prompting}. An example prompt is provided in \cref{fig:pull-figure}.

We offer empirical results on arithmetic, commonsense, and symbolic reasoning benchmarks, demonstrating that chain-of-thought prompting surpasses standard prompting, sometimes by a substantial margin. \cref{fig:pull-bar-chart} displays one such example—on the GSM8K dataset of math word problems \citep{cobbe2021training}, chain-of-thought prompting with \palm{} 540B significantly outperforms standard prompting and sets a new state-of-the-art. A purely prompting-based method is valuable because it does not demand a large labeled training set and allows a single model checkpoint to address multiple tasks without losing generality. This work highlights how large language models can acquire knowledge from a few examples expressed in natural language concerning the task, as opposed to relying solely on discovering patterns from a large training corpus.

\section{Chain-of-Thought Prompting} When tackling a complex reasoning problem, such as a multi-step math word question, people often reflect on their own thought process. Typically, the problem is broken down into smaller intermediate steps that are solved sequentially before arriving at the final solution: \textit{``After Jane gives 2 flowers to her mom, she has 10 $\ldots$ then after she gives 3 to her dad, she will have 7 $\ldots$ so the answer is 7.''} This paper aims to enable language models to produce a comparable \textit{chain of thought}—a logically connected sequence of intermediate reasoning steps culminating in the problem’s final answer. We demonstrate that large enough language models can generate such chains of thought when provided with examples of chain-of-thought reasoning in few-shot prompting.

\cref{fig:pull-figure} illustrates an instance where a model creates a chain of thought to solve a math word problem it would have otherwise answered incorrectly. Here, the chain of thought resembles a solution and can be understood as one, but we deliberately term it a chain of thought to emphasize that it simulates a step-by-step reasoning process leading to the answer (and, notably, solutions or explanations generally follow the final answer \citep[][\textit{inter alia}]{narang2020wt5,wiegreffe2021reframing,lampinen2022can}).

Chain-of-thought prompting offers multiple appealing advantages as a method to enhance reasoning capabilities in language models.

\begin{enumerate}[topsep=1pt,itemsep=0ex]
    \item Firstly, chain of thought, fundamentally, enables models to break down multi-step problems into sequential intermediate steps, allowing additional computation to be dedicated to problems that demand more reasoning stages.
    \item Secondly, a chain of thought offers an interpretable insight into the model’s behavior, indicating how it may have reached a specific answer and providing chances to diagnose where the reasoning process might have faltered (although completely defining the computations a model uses to arrive at an answer remains unresolved).
    \item Thirdly, chain-of-thought reasoning can be applied to tasks such as math word problems, commonsense reasoning, and symbolic manipulation, and is potentially relevant (at least theoretically) to any task humans solve through language.
    \item Lastly, chain-of-thought reasoning can be effortlessly elicited from sufficiently large, off-the-shelf language models simply by incorporating chain of thought sequence examples within the few-shot prompting exemplars.
\end{enumerate}

Through empirical studies, we will demonstrate the effectiveness of chain-of-thought prompting in arithmetic reasoning (\cref{sec:arithmetic-reasoning}), commonsense reasoning (\cref{sec:commonsense-reasoning}), and symbolic reasoning (\cref{sec:symbolic-reasoning}).

\section{Arithmetic Reasoning}\label{sec:arithmetic-reasoning}
We start by examining math word problems similar to those illustrated in \cref{fig:pull-figure}, which assess the arithmetic reasoning capabilities of language models. While straightforward for humans, arithmetic reasoning remains a challenge for language models \citep[][\textit{inter alia}]{hendrycks2021measuring,patel-etal-2021-nlp}. Remarkably, when paired with the 540B parameter language model, chain-of-thought prompting achieves performance comparable to task-specific finetuned models on several benchmarks, even setting a new state-of-the-art on the difficult GSM8K dataset \citep{cobbe2021training}.

\subsection{Experimental Setup}

We investigate chain-of-thought prompting across various language models and multiple evaluation sets.

\textbf{Benchmarks.}
We evaluate the following five math word problem benchmarks:
\textbf{(1)} the \textbf{GSM8K} math word problems dataset \citep{cobbe2021training}, 
\textbf{(2)} the \textbf{SVAMP} dataset featuring math problems with varied structures \citep{patel-etal-2021-nlp},
\textbf{(3)} the \textbf{ASDiv} collection of diverse math word problems \citep{miao-etal-2020-diverse},
\textbf{(4)} the \textbf{AQuA} set of algebraic word problems, and
\textbf{(5)} the \textbf{MAWPS} benchmark \citep{koncel-kedziorski-etal-2016-mawps}. Sample problems appear in Appendix \cref{tab:math-datasets}.

\textbf{Standard prompting.} For the baseline, we utilize standard few-shot prompting as introduced by \citet{brown2020language}, where a language model is provided with in-context exemplars of input-output pairs before generating a response for a test example. Exemplars are formatted as question-and-answer pairs. The model outputs the final answer directly, as depicted in \cref{fig:pull-figure} (left).

\textbf{Chain-of-thought prompting.} Our method involves enhancing each example in few-shot prompting by including a chain of thought linked to the corresponding answer, as depicted in \cref{fig:pull-figure} (right). Since most datasets only provide an evaluation split, we manually created a set of eight few-shot examples with chains of thought for prompting—\cref{fig:pull-figure} (right) presents one such chain of thought example, and the complete collection of examples is available in Appendix \cref{tab:appendix-mwp-prompts}. (These specific examples were not optimized through prompt engineering; robustness is examined in \cref{subsec:robustness} and \cref{subsec:faq-prompt-engineering}.) To assess whether chain-of-thought prompting in this format effectively triggers accurate reasoning across various math word problems, we utilized this single set of eight chain-of-thought examples for all benchmarks except AQuA, which is multiple-choice rather than free response. For AQuA, we employed four examples and their solutions from the training set, as detailed in Appendix \cref{tab:appendix-aqua-prompt}.

\textbf{Language models.} We test five large language models. The first is \textbf{GPT-3} \citep{brown2020language}, using the variants text-ada-001, text-babbage-001, text-curie-001, and text-davinci-002, which likely correspond to InstructGPT models with 350M, 1.3B, 6.7B, and 175B parameters respectively \citep{ouyang2022training}. The second is \textbf{\lamda{}} \citep{thoppilan2022lamda}, available in 422M, 2B, 8B, 68B, and 137B parameter sizes. The third is \textbf{\palm{}}, with models of 8B, 62B, and 540B parameters. The fourth is \textbf{UL2 20B} \citep{tay2022unifying}, and the fifth is \textbf{Codex} \citep[][code-davinci-002 in the OpenAI API]{chen2021evaluating}. We generate outputs from these models via greedy decoding (although subsequent research indicates chain-of-thought prompting can be enhanced by selecting the most frequent final answer across multiple sampled outputs \citep{wang2022self}). For \lamda{}, we present results averaged over five random seeds, each using a differently randomized order of examples. Since \lamda{} experiments showed minimal variation across seeds, to conserve computational resources, we report results for a single example order for all other models.

\subsection{Results}  
The most notable findings regarding chain-of-thought prompting are presented in \cref{fig:main-math}, while the comprehensive outputs for each model collection, model size, and benchmark are detailed in \cref{tab:all-lm-math} in the Appendix. There are three primary insights. First, \cref{fig:main-math} illustrates that chain-of-thought prompting emerges as a capability with increased model scale \citep{wei2022emergent}. Specifically, this technique does not enhance performance in smaller models and only shows improvements when applied to models with approximately 100B parameters. Our qualitative analysis revealed that smaller-scale models generated coherent but illogical chains of thought, resulting in worse performance compared to standard prompting.

\input{main_fables/main-math}  
Second, the benefits of chain-of-thought prompting are more pronounced on complex problems. For example, on GSM8K—the dataset with the lowest baseline accuracy—the largest GPT and \palm{} models more than doubled their performance. Conversely, on SingleOp, the simplest subset of MAWPS that requires only one step to solve, improvements were minimal or even negative (see Appendix \cref{tab:mawps-subsets}).

Third, when used with GPT-3 175B and \palm{} 540B, chain-of-thought prompting competes favorably with previous state-of-the-art methods, which typically rely on fine-tuning task-specific models using labeled training data. \cref{fig:main-math} demonstrates that \palm{} 540B employing chain-of-thought prompting sets new state-of-the-art results on GSM8K, SVAMP, and MAWPS (noting that standard prompting already surpassed previous best results on SVAMP). For the remaining datasets, AQuA and ASDiv, \palm{} with chain-of-thought prompting achieves performance within 2\% of the state of the art (Appendix \cref{tab:all-lm-math}).

To gain a deeper understanding of why chain-of-thought prompting is effective, we conducted a manual review of the reasoning chains generated by \lamda{} 137B for GSM8K. Among 50 randomly selected instances where the model produced the correct final answer, every chain of thought was both logically and mathematically sound, except for two cases that happened to reach the correct answer by chance (refer to \cref{subsec:correct-chain-of-thought} and \cref{tab:appendix-gsm8k-correct-examples} for examples of accurate model-generated reasoning chains). Additionally, we examined 50 random cases in which the model output an incorrect answer. Our analysis revealed that 46\% of these reasoning chains were nearly correct, containing only minor errors such as calculation mistakes, symbol mapping errors, or a single omitted reasoning step, while the remaining 54\% had significant flaws related to semantic understanding or coherence (see \cref{subsec:incorrect-chain-of-thought-analysis}). 

To shed some light on why increasing model scale enhances chain-of-thought reasoning capabilities, we performed a comparable error analysis on \palm{} 62B, examining which errors were rectified upon scaling up to \palm{} 540B. The findings indicate that scaling \palm{} to 540B resolves a substantial portion of the errors involving missing single steps and semantic misunderstandings present in the 62B model (see \cref{subsec:why-scale-helps}).

\subsection{Ablation Study}\label{subsec:math-ablation}

The demonstrated advantages of chain-of-thought prompting naturally prompt the question of whether similar performance gains can be achieved through alternative prompting methods. \cref{fig:bar_ablation} presents an ablation study featuring three variants of chain-of-thought prompting, outlined below.

\textbf{Equation only.} One potential reason chain-of-thought prompting is beneficial is that it generates the mathematical equation needed for evaluation. Therefore, we experimented with a variant where the model is prompted to produce only the mathematical equation prior to providing the final answer. \cref{fig:bar_ablation} reveals that equation-only prompting offers limited improvement for GSM8K, suggesting that the semantics of GSM8K questions are too complex to be directly converted into an equation without the intermediate natural language reasoning steps characteristic of chain of thought. However, for datasets consisting of problems requiring only one or two steps, we observe that equation-only prompting does enhance performance, since the equation can be straightforwardly derived from the question (see Appendix \cref{tab:ablations-arithmetic}).

\input{main_fables/ablation-bar}
\textbf{Variable compute only.} One hypothesis is that chain of thought enables the model to allocate more computation (i.e., intermediate tokens) on more challenging problems. To separate the influence of variable computation from the reasoning process in chain-of-thought, we evaluate a setting where the model is prompted to generate only a sequence of dots ($\ldots$) matching the number of characters in the equation required to solve the problem. This version performs comparably to the baseline, indicating that variable computation alone does not explain the effectiveness of chain-of-thought prompting, and that articulating intermediate steps in natural language provides additional value.

\textbf{Chain of thought after answer.} Another possible advantage of chain-of-thought prompting is that it might help the model better retrieve relevant knowledge obtained during pretraining. To investigate this, we test an alternative setup where the chain-of-thought prompt is provided only after the answer, which isolates whether the model relies on the generated chain of thought to produce the final response. This variant performs similarly to the baseline, suggesting that the sequential reasoning represented by the chain of thought contributes benefits beyond merely triggering knowledge activation.

\input{main_fables/main-robustness}
\subsection{Robustness of Chain of Thought}\label{subsec:robustness}

The sensitivity to the choice of exemplars is a crucial factor in prompting methods—for example, changing the order of few-shot exemplars can lead to GPT-3’s accuracy on SST-2 varying from nearly random (54.3\%) to nearly state-of-the-art (93.4\%) \citep{zhao2021calibrate}. In this last subsection, we analyze the robustness of chain of thought prompts authored by different annotators. Beyond the results already presented, which used chains of thought written by Annotator A, two other co-authors (Annotators B and C) independently created chains of thought for the same few-shot examples (see \cref{sec:appendix-alternate-annotators}). Annotator A also authored an additional, more concise chain of thought following the style of solutions in \citet{cobbe2021training}.\footnote{For example, while the original chain of thought uses multiple short sentences (\textit{``There were originally 9 computers. For each of 4 days, 5 more computers were added. So 5 * 4 = 20 computers were added. 9 + 20 is 29.''}), the concise version would state \textit{``5 * 4 = 20 new computers were added. So there are 9 + 20 = 29 new computers in the server room now''}.}

\cref{fig:bar-robustness-analysis} presents these findings for \lamda{} 137B on the GSM8K and MAWPS datasets (additional ablation results for other datasets can be found in Appendix \cref{tab:ablations-arithmetic} / \cref{tab:ablations-commonsense-symbolic}). Despite the expected variability among different chain of thought annotations, as is typical when using exemplar-based prompting \citep{le-scao-rush-2021-many,reynolds2021prompt,zhao2021calibrate}, every set of chain of thought prompts significantly outperforms the standard baseline. This indicates that effective use of chain of thought prompting does not rely on a specific linguistic style.

To verify that chain-of-thought prompting remains effective with other exemplar sets, we conducted experiments using three sets of eight exemplars randomly selected from the GSM8K training set, an independent source whose examples already include reasoning steps resembling chain of thought.\footnote{We selected examples with $\leq 60$ tokens to fit within our input context window and restricted examples to $\leq 2$ steps to ensure a fair comparison with the eight exemplars we composed.} \cref{fig:bar-robustness-analysis} demonstrates that these randomly sampled prompts performed on par with our manually created exemplars and also significantly outperformed standard prompting.

Beyond robustness to different annotators, independently authored chains of thought, varied exemplars, and multiple language models, we also observe that chain-of-thought prompting for arithmetic reasoning is stable across different exemplar orders and different numbers of exemplars (refer to \cref{subsec:faq-prompt-engineering}).

\section{Commonsense Reasoning}\label{sec:commonsense-reasoning}

While chain of thought is especially effective for math word problems, its language-based approach extends naturally to a wide range of commonsense reasoning tasks. These tasks involve reasoning about physical and social interactions based on general background knowledge. Commonsense reasoning is crucial for navigating the world and remains beyond the capabilities of current natural language understanding systems \citep{talmor2022commonsenseqa}.

\textbf{Benchmarks.}  
We evaluate five datasets that encompass a wide variety of commonsense reasoning challenges. The well-known \textbf{CSQA} \citep{talmor-etal-2019-commonsenseqa} involves answering commonsense questions about the world, often requiring complex semantic understanding and prior knowledge.  
\textbf{StrategyQA} \citep{geva-etal-2021-aristotle} demands that models deduce a multi-step reasoning strategy to solve the questions. We also include two specialized evaluation datasets from the BIG-bench collection \citep{bigbench}: \textbf{Date} Understanding, which entails inferring dates based on context, and \textbf{Sports} Understanding, which tests the plausibility of sentences related to sports. Lastly, the \textbf{SayCan} dataset \citep{ahn2022can} focuses on translating natural language instructions into a sequence of discrete robot actions.  
\cref{fig:dataset-examples} provides sample instances with chain of thought annotations for each dataset.

\textbf{Prompts.}  
Our experimental procedure follows that of the previous section. For CSQA and StrategyQA, we randomly chose examples from the training splits and manually created chains of thought to serve as few-shot exemplars. Since the two BIG-bench tasks lack training sets, we selected the first ten examples from their evaluation sets as few-shot exemplars and report results on the remaining evaluation examples. For SayCan, we utilized six examples from the training data used in \citet{ahn2022can} and similarly crafted chains of thought manually.

\textbf{Results.}  
\cref{fig:commonsense-results} showcases these findings for \palm{} (comprehensive results for \lamda{}, GPT-3, and various model sizes are provided in \cref{tab:all-lm-commonsense}).  
Across all tasks, increasing the model size enhanced performance with standard prompting; applying chain-of-thought prompting yielded additional improvements, with the most significant gains observed for \palm{} 540B. Using chain-of-thought prompting, \palm{} 540B demonstrated strong results compared to baselines, surpassing the previous state of the art on StrategyQA (75.6\% versus 69.4\%) and outperforming an unaided sports fan on sports understanding (95.4\% versus 84\%).  
These outcomes illustrate that chain-of-thought prompting can also boost performance on tasks that require a diverse set of commonsense reasoning skills (although the improvement was minimal on CSQA).

\input{main_fables/main-commonsense}

\input{main_fables/main-symbolic}  
\section{Symbolic Reasoning}\label{sec:symbolic-reasoning}  
Our final experimental investigation focuses on symbolic reasoning, which is straightforward for humans but may pose difficulties for language models. We demonstrate that chain-of-thought prompting not only enables language models to handle symbolic reasoning tasks that are difficult under standard prompting conditions but also supports length generalization to inference-time inputs longer than those presented in the few-shot examples.

\paragraph{Tasks.}  
We consider the following two synthetic tasks.

\begin{itemize}[leftmargin=*,topsep=0pt]  
    \itemsep0em \item \textbf{Last letter concatenation.} This task requires the model to join the last letters of each word in a name (e.g., \textit{``Amy Brown''} $\rightarrow$ \textit{``yn''}).  
    It represents a more difficult variant of first letter concatenation, which language models can already accomplish without chain-of-thought prompting.\footnote{We evaluated 10 common names using GPT-3 \texttt{davinci} and it correctly answered all but one.} Full names are generated by randomly combining names from the top one-thousand first and last names based on name census data ({\small{\url{https://namecensus.com/}}}).  
    \item \textbf{Coin flip.} This task asks the model to determine if a coin remains heads up after individuals either flip or do not flip the coin (e.g., \textit{``A coin is heads up. Phoebe flips the coin. Osvaldo does not flip the coin. Is the coin still heads up?''} $\rightarrow$ \textit{``no''}).  
\end{itemize}

Since the design of these symbolic reasoning tasks is clearly defined, for each task we create an \textit{in-domain} test set where the examples have the same number of steps as those in the training or few-shot exemplars, as well as an \textit{out-of-domain} (OOD) test set where the evaluation examples contain more steps than those present in the exemplars. In the case of last letter concatenation, the model is exposed only to exemplars consisting of two-word names and is then tested on last letter concatenation for names with 3 and 4 words.\footnote{For names longer than two words, we concatenate multiple first and last names together.} A similar approach is used for varying the number of potential flips in the coin flip task. Our experimental framework employs the same methods and models as those described in the previous two sections. We once again manually construct chains of thought for the few-shot exemplars specific to each task, as provided in \cref{fig:dataset-examples}.

\paragraph{Results.} 
The outcomes of both in-domain and OOD evaluations for \palm{} are presented in \cref{fig:symbolic-main}, with corresponding results for \lamda{} included in Appendix \cref{tab:all-lm-symbolic}. Using \palm{} 540B, chain-of-thought prompting achieves nearly 100\% success rates (it should be noted that standard prompting already succeeds on the coin flip task with \palm{} 540B, unlike with \lamda{} 137B). It is important to emphasize that these in-domain evaluations are considered "toy tasks" because the perfect solution structures are already demonstrated by the chains of thought in the few-shot exemplars; the model merely needs to replicate these steps with the new symbols in the test example. Nevertheless, smaller models still fail— the capacity to carry out abstract manipulations on previously unseen symbols in these three tasks only emerges at the scale of models with approximately 100B parameters.

Regarding the OOD evaluations, standard prompting fails on both tasks. However, with chain-of-thought prompting, language models exhibit upward scaling trends (although performance remains lower than in the in-domain scenario). Thus, chain-of-thought prompting enables length generalization beyond previously seen reasoning chains for sufficiently large language models.

\section{Discussion}\label{sec:discussion}

We have investigated chain-of-thought prompting as a straightforward method to invoke multi-step reasoning in large language models. Initially, we observed that chain-of-thought prompting significantly enhances performance in arithmetic reasoning, producing improvements far greater than those seen in ablation studies and demonstrating robustness across different annotators, exemplars, and language models (\cref{sec:arithmetic-reasoning}). Subsequently, experiments on commonsense reasoning highlighted how the linguistic structure of chain-of-thought reasoning renders it broadly applicable (\cref{sec:commonsense-reasoning}). Finally, we demonstrated that for symbolic reasoning, chain-of-thought prompting supports OOD generalization to longer sequence lengths (\cref{sec:symbolic-reasoning}). In all cases, chain-of-thought reasoning is induced merely by prompting an off-the-shelf language model, with no fine-tuning performed during this study.

The emergence of chain-of-thought reasoning as a function of model scale has been a consistent observation \citep{wei2022emergent}. For numerous reasoning tasks where standard prompting exhibits flat scaling, chain-of-thought prompting results in substantially rising scaling curves. It appears that chain-of-thought prompting broadens the range of tasks large language models can solve effectively—in other words, our findings emphasize that standard prompting merely sets a lower bound on the capabilities of large language models. This insight likely raises more questions than it resolves—for example, how much further might reasoning abilities improve with increased model scale? What additional prompting techniques could extend the scope of tasks that language models can address?

Regarding limitations, we first note that although chain-of-thought mimics human reasoning processes, this does not establish that the neural network is genuinely ‘‘reasoning,’’ which remains an open question. Second, while the manual effort to augment exemplars with chains of thought is minimal in few-shot settings, such annotation costs could become prohibitive for fine-tuning—although this might be alleviated through synthetic data generation or zero-shot generalization. \orange{Third, there is no guarantee that the reasoning paths are correct, which can yield both correct and incorrect answers; enhancing the factual accuracy of language model generations is an important avenue for future research \citep[][\textit{inter alia}]{rashkin2021measuring,ye2022unreliability,wiegreffe2021reframing}.} Finally, since chain-of-thought reasoning only emerges in large-scale models, deploying it in real-world applications can be costly; future work might investigate methods to induce reasoning capabilities in smaller models.

\section{Related Work} This study draws inspiration from various research fields, which we discuss in more detail in an extended related work section (\cref{sec:extended-related-work}). Here, we highlight two key areas and the relevant literature most closely connected to our work.

The first pertinent area involves using intermediate steps to address reasoning challenges. \cite{ling-etal-2017-program} introduce the concept of employing natural language rationales to solve math word problems through a sequence of intermediate stages. Their approach stands in stark contrast to prior work that utilizes formal languages for reasoning \citep{roy-2016-reasoning, chiang-chen-2019-semantically, amini-etal-2019-mathqa, chen2019neural}. Building on \citet{ling-etal-2017-program}, \citet{cobbe2021training} develop a larger dataset and leverage it to finetune a pretrained language model rather than training from scratch. In the area of program synthesis, \cite{nye2021show} exploit language models to predict the final outputs of Python programs by first predicting intermediate computational steps line-by-line, demonstrating that this stepwise prediction method surpasses directly forecasting final outputs.

Additionally, this paper is closely connected to a substantial body of recent work on prompting. Following the popularization of few-shot prompting by \citet{brown2020language}, several general strategies have been proposed to enhance model prompting capabilities, including automatically learned prompts \citep{lester-etal-2021-power} and providing models with explicit task instructions \citep{wei2021finetuned,sanh2021multitask,ouyang2022training}. While these techniques focus on improving or enriching the input segment of the prompt (such as instructions appended before inputs), our approach takes a complementary path by enhancing the outputs of language models with a chain of thought.

\section{Conclusions}
We have investigated chain-of-thought prompting as a straightforward and widely applicable technique to improve reasoning abilities in language models. Through experiments covering arithmetic, symbolic, and commonsense reasoning tasks, we observe that chain-of-thought reasoning emerges with increasing model scale, enabling sufficiently large language models to solve reasoning problems that otherwise show flat scaling performance. Expanding the range of reasoning tasks addressable by language models may encourage further exploration of language-driven methods for reasoning.

\end{document}